\documentclass[11pt]{article}
\usepackage[a4paper,margin=1in]{geometry}
\usepackage{amsmath,amssymb,amsthm,mathtools}
\usepackage[hidelinks]{hyperref}
\usepackage{booktabs}

\newtheorem{theorem}{Theorem}
\newtheorem{proposition}[theorem]{Proposition}
\newtheorem{lemma}[theorem]{Lemma}
\newtheorem{conjecture}[theorem]{Conjecture}
\newcommand{\Q}{\mathbb Q}
\newcommand{\Z}{\mathbb Z}
\newcommand{\Tr}{\operatorname{Tr}}

\title{Dimension-Six Twisted Convolution as Arithmetic Boundary Fusion\\
of a Two-Base Lens-Space Beta Integral}
\author{Hainan Zhao\\
\small Independent Researcher\\
\small\texttt{hainzhao@gmail.com}}
\date{29 July 2026}

\begin{document}
\maketitle

\begin{abstract}
We reformulate the remaining dimension-six case of the formal Twisted
Convolution Conjecture of Appleby--Flammia--Kopp as one boundary theorem
for a two-base general modular quantum-dilogarithm packet.  The
dimension-six principal-ghost quotient is identified with the
two-gamma kernel in the degenerate Sarkissian--Spiridonov lens-space
beta transform.  In the upper half-plane its natural bases are
$e^{2\pi i\tau}$ and $e^{2\pi iA_6\tau}$ and must not be fused.
At the real-multiplication fixed point
$\beta_6+\beta_6^{-1}=5$ they coalesce arithmetically and give the
previously isolated well-poised ${}_2\psi_2$ packet at argument $-q$.
The identical calculation in the proved dimension-five case lands
instead on the summable $+q$ locus, while the independent
dimension-four even-wrap control returns to $-q$.  We verify the
interior identity as a meromorphic spectral identity, specialize the
published beta
evaluation exactly, and show that the undeformed endpoint contour is
not absolutely convergent because its quadratic decay cancels.  We
define the finite part by an interior tilted contour and prove
tilt-independence by Cauchy's theorem in the pole-free strip.  An
exact all-characteristic calculation gives
$A_6\equiv I\pmod6$, $\psi^2(A_6)=-1$, and equality of the Kopp and AFK
multipliers in all thirty-six cases.  Consequently one explicitly
stated analytic fusion-continuity conjecture implies both formal
dimension-six shifts.  We do not claim that conjecture here.
\end{abstract}

\section{Statement and proof status}

Put
\[
 \beta=\frac{5+\sqrt{21}}2,\qquad
 A=\begin{pmatrix}115&-24\\24&-5\end{pmatrix}.
\tag{1}
\]
Then
\[
 \beta^2-5\beta+1=0,\qquad A\beta=\beta,\qquad
 \beta+\beta^{-1}=5.
\tag{2}
\]
The first equality is the arithmetic trace relation; throughout this
paper it is part of the boundary condition, not a generic removable
limit.

The finite part of the dimension-six calculation is already exact.
There are two formal shifts, the reconstructed matrix has trace one,
and all $225$ rank-two minors vanish after imposing the certified
algebraic packet.  What is missing is the analytic identification of
that packet with the convention-matched AFK cocycle values.  The
purpose of this note is to make the remaining analytic statement
independent of the finite certificates.

Parametrize the attracting half of the $A$-axis by
\[
 \gamma(s)=\frac52+
 \frac{\sqrt{21}}2\frac{1-s^2}{1+s^2}
 +i\frac{\sqrt{21}s}{1+s^2},\qquad s>0.
\]
Then $\gamma(s)\to\beta$ as $s\downarrow0$.
Kopp's Definition~4.7 chooses the positive stabilizer generator by this
attracting direction and observes that it translates the same geodesic
toward $\beta$~\cite{Kopp}.  Definition~4.9 defines the RM stable value
by evaluation of the meromorphic continuation at the fixed point, and
Proposition~7.20 realizes it as a limit from the upper half-plane.

Let $\mathcal P_j(\tau)$, $j=0,1,2$, denote the primitive ray-class
logarithms of the spectral periodization in the norm-$37$ Frobenius
order, with logarithm branches continued from the two-base chamber
along $\gamma$, and put
\[
 \mathcal R_1(\tau)=
 \mathcal P_0(\tau)+\zeta_6\mathcal P_1(\tau)
 +\zeta_6^2\mathcal P_2(\tau).
\]

For $\tau$ in the interior convergence chamber, put
$S_\tau=\{y:0<\Re y<\Re Q_\tau\}$.  If
\[
 C_h=\{c+h(t)+it:t\in\mathbb R\}\subset S_\tau
\]
is a $C^1$ graph homotopic to $c+i\mathbb R$ through graphs on which
the two-base kernel has one common integrable majorant, define its
\emph{tilted value} by integration over $C_h$, followed by the exact
helical periodization.

\begin{proposition}[Tilt independence]\label{prop:tilt}
The interior tilted value is independent of the admissible graph
$C_h$.
\end{proposition}

\begin{proof}
Truncate two admissible graphs at imaginary heights $\pm T$ and join
their endpoints inside $S_\tau$.  The kernel is holomorphic in the
region between them because its pole cones lie on opposite sides of
the strip.  Cauchy's theorem makes the integral around the resulting
boundary zero.  The common interior exponential majorant makes both
joining integrals tend to zero as $T\to\infty$.  The two graph
integrals are equal.  Absolute local uniform convergence permits the
same argument after helical periodization.
\end{proof}

Thus $\mathcal P_j(\tau)$ may equivalently be defined by any admissible
tilt, and
\[
 \operatorname{FP}_{\beta}\mathcal P_j
 :=\lim_{s\downarrow0}\mathcal P_j(\gamma(s)),
\]
if it exists, has no remaining contour-choice ambiguity.

For each \(j\), let \(U_j^{\rm fus}\) be the boundary value obtained by
first taking the meromorphic spectral periodization with its certified
level-\(24\) label, then imposing the fixed-point relation
\(\widetilde q_M=q_M\), and evaluating the resulting expression through
the \(24\)-factor continuation rather than the collapsed bilateral
series.  The logarithm is the continuation of the branch used for
\(\mathcal P_j\).  Define
\[
 \operatorname{Fus}_5\mathcal R_1
 :=\log U_0^{\rm fus}
 +\zeta_6\log U_1^{\rm fus}
 +\zeta_6^2\log U_2^{\rm fus}.
\]
The subscript \(5\) records the trace equation
\(\beta+\beta^{-1}=5\), not a dimension.

\begin{conjecture}[Minimal fusion-continuity, MFC$_6$]
\label{conj:fusion}
The scalar $\mathcal R_1(\gamma(s))$ has a finite limit as
$s\downarrow0$, and meromorphic spectral periodization commutes in this
primitive component with trace-five base fusion, preserving the
norm-$37$ lens/Frobenius label:
\[
 \lim_{s\downarrow0}\mathcal R_1(\gamma(s))
 =\operatorname{Fus}_5\mathcal R_1.
\]
\end{conjecture}

No uniform convergence, H\"older estimate, or two-sided continuity is
assumed.  Since $A$ translates the same oriented geodesic toward its
attracting fixed point, existence of the limit makes it $A$-invariant;
invariance is not an additional hypothesis.  The exceptional zero
frequency is also not part of MFC$_6$: it is supplied separately by the
exact AFK endpoint $-4\sqrt7$.  Its tilted Fresnel/Abel value gives the
reciprocal pair $-2\sqrt7\pm3\sqrt3$, whose trace is $-4\sqrt7$.

\begin{theorem}[Analytic-to-Stark bridge]\label{thm:stark}
Assume Conjecture~\ref{conj:fusion}.  Let
$G=\operatorname{Cl}_{(6)\infty_2}(\Q(\sqrt{21}))
=\langle g\rangle\simeq C_6$, let
$\chi_1(g)=\zeta_6$, and let $r_j$ be the logarithms of the certified
algebraic ray-unit orbit in the arithmetic Frobenius order.  Then
\[
 L'_S(0,\chi_1)
 =r_0+\zeta_6r_1+\zeta_6^2r_2.
\tag{3}
\]
Thus a flow-invariant analytic regularity statement at the
$A_6$ fixed point implies the convention-fixed rank-one Stark
identity over $\Q(\sqrt{21})$.
\end{theorem}

\begin{theorem}[Conditional dimension-six closure]\label{thm:conditional}
Assume Conjecture~\ref{conj:fusion}.  Then $0$ and $1$ are shifts for
every admissible dimension-six tuple in the sense of
AFK~\cite[Definition 1.34 and equation (1.49)]{AFK}.
\end{theorem}

The implication in Theorem~\ref{thm:conditional} is proved in
Section~\ref{sec:conditional}.  Conjecture~\ref{conj:fusion} remains
open.  In particular, the theorem is not an unconditional proof of the
dimension-six TCC.

\section{The honest two-base packet}\label{sec:two-base}

For the general modular quantum dilogarithm of
Sarkissian--Spiridonov~\cite{SS}, use
\[
 (p,k,r,s)=(-115,24,5,24),\qquad pr+ks=1.
\tag{4}
\]
Their equations (5)--(8) give the two bases
\[
 q_M=e^{2\pi i\tau},\qquad
 \widetilde q_M=e^{2\pi iA\tau}.
\tag{5}
\]
They are distinct in the upper half-plane.  Their equation (15)
factors the same general modular function into $24$ standard Faddeev
factors.  If $\rho=24\tau-5$, the bases in each standard factor are
\[
 q_S=e^{2\pi i\rho},\qquad
 \widetilde q_S=e^{-2\pi i/\rho}.
\tag{6}
\]
Equations (4) and (5) are different modular pairs away from the
boundary.  Conflating them gives the false equal-base interior model.

\begin{proposition}[Interior factorization]\label{prop:interior}
In the chamber in which both product bases have modulus less than one,
the direct general modular gamma in (4) equals its $24$-factor
Faddeev continuation.  Periodizing the degenerate two-gamma kernel
over the helical quotient
\[
 \bigl(\mathbb R\times\Z/24\Z\bigr)/
 \langle(\omega_1-\omega_2,6)\rangle
\tag{7}
\]
gives three absolutely and locally uniformly convergent bibasic
bilateral classes for each level-six frequency.  The resulting
identity is valid as a meromorphic spectral identity.
\end{proposition}

\begin{proof}
The factorization is equation (15) of~\cite{SS} specialized by (3).
For a frequency $(a,b)\bmod6$ and alias $z\in\Z$, direct reduction gives
\[
 N_z=a+2-6z,\qquad \ell_z=b-6z,
\tag{8}
\]
and
\[
 \frac{\alpha_z}{4\tau-1}
 =\frac{4b-5a}{3}+2z.
\tag{9}
\]
Thus $-5(N_z-2)\equiv a$ and $\ell_z\equiv b\pmod6$.
Three alias steps are one helical translation at level $24$.
The functional-equation ratio is bibasic and has geometric tails in
the stated chamber.  Finite heads and the bound
\[
 \frac{|a|^{M+1}}
 {(M+1)(1-|a|)(1-|q|)}
\tag{10}
\]
for the omitted logarithmic $q$-Pochhammer tail prove local uniform
convergence.  This justifies the spectral reindexing.  The exact
integer ledger and Arb enclosures are reproduced by
\path{scripts/dimension_six_two_base_lens.py} and
\path{scripts/dimension_six_interior_factorization_audit.py}.
\end{proof}

At $\tau=\beta$, both pairs fuse, but for different reasons:
$A\beta=\beta$ for (4), while
$\beta+\beta^{-1}=5$ makes the bases in (5) equal.  Exact substitution
in the term ratio gives
\[
 -q\,\frac{(1-x)(1+w^{-1}x)}
 {(1+qx)(1-qw^{-1}x)}.
\tag{11}
\]
This is the well-poised ${}_2\psi_2$ packet at argument $-q$ found in
the finite Zak descent.  Equation (11) is a boundary identity only.

\section{The dimension-five control}

For dimension five,
\[
 \beta_5=2+\sqrt3,\qquad
 A_5=\begin{pmatrix}56&-15\\15&-4\end{pmatrix},\qquad
 \beta_5+\beta_5^{-1}=4.
\tag{12}
\]
Running the same two-base construction gives at fusion
\[
 q\,\frac{(1-x)(1-w^{-1}x)}
 {(1-qx)(1-qw^{-1}x)}.
\tag{13}
\]

\begin{proposition}[Summability dichotomy]\label{prop:dichotomy}
The proved dimension-five packet lies on the positive-argument
closed ${}_2\psi_2$ locus (13).  The dimension-six packet lies on its
sign-reflected neighbor (11).
\end{proposition}

\begin{proof}
The sign follows from the exact even-wrap and lens-label reductions.
For dimension five the direct two-base function and its factorization
are enclosed at three interior points, all six bibasic classes
converge with rigorous tails, and the boundary double-sine square
encloses the independently proved algebraic root
\[
 3.8908617139430792553376439596\ldots.
\]
The executable proof is
\path{scripts/dimension_five_two_base_calibration.py}.  The
dimension-four control independently fixes lens level $8$ and phase
level $16$.  Exact exponent reduction also gives
\[
 -q\,\frac{(1-x)(1-ix)}{(1+qx)(1+iqx)},
\]
so dimension four, like dimension six, lies at argument $-q$; see
\path{scripts/dimension_four_two_base_calibration.py}.
The exceptional dimension-six zero mode is independently normalized by
\[
 \nu_{\rm aux}+\nu_{\rm aux}^{-1}=-4\sqrt7.
\]
Thus the complete calibration triangle is
\[
 d=4:-q,\qquad d=5:+q,\qquad
 d=6\text{ zero mode}:-4\sqrt7.
\]
\end{proof}

This control is decisive: the minus sign in (11) is not a dispensable
normalization, and replacing it by the summable plus sign would erase
the dimension-six obstruction.

\section{Exact beta evaluation and its residue}

Sarkissian--Spiridonov's unnumbered beta-evaluation theorem immediately
preceding their equation (58), followed by their degeneration (66),
applies meromorphically after the specialization
\[
 \omega_1=24\beta-5=\beta^3,\quad
 \omega_2=1,\quad Q=\omega_1+\omega_2,\quad
 g=Q,\quad \ell=0.
\tag{14}
\]
The continuous phase coefficient is
\[
 p-k(1-s)=-115-24(1-24)=437.
\tag{15}
\]
For the Fourier variables $(\alpha,N)$, equation (66) becomes
\[
\begin{split}
 &\sum_{m\bmod24}\int
 e^{\pi i m\,437(2N-4)/24}
 e^{\pi i\alpha(2y-Q)/(24\omega_1)}
 \Gamma_M(y,m)\Gamma_M(Q-y,-m)
 \frac{dy}{i\sqrt{\omega_1}}\\
 &\qquad =
 24\Gamma_M(-\alpha,4-N)
 \Gamma_M(\alpha,N)\Gamma_M(Q,0),
\end{split}
\tag{16}
\]
as a meromorphic Fourier identity.

\begin{proposition}[Conservation of orientation]\label{prop:residue}
The right side of (15) does not reduce, by the standard double-sine
reflection and period-shift relations, to an unoriented norm identity.
The irreducible oriented residue is
\[
 \Gamma_M(-\alpha,4-N)\Gamma_M(\alpha,N).
\tag{17}
\]
\end{proposition}

\begin{proof}
The reflection partner of $\Gamma_M(\alpha,N)$ is
$\Gamma_M(Q-\alpha,4-N)$, not
$\Gamma_M(-\alpha,4-N)$.  Shifting by $Q$ changes the discrete label by
$r-1=4$, which is not zero modulo $24$.  Canonical reduction therefore
leaves (16).  The full parameter and relation-basis audit is
\path{scripts/dimension_six_ss_evaluation_audit.py}.
\end{proof}

Thus (15) is a genuine integral-transform identity but not, on its own,
the new finite multiplicative relation required to identify the
dimension-six algebraic packet.

\section{Why the endpoint contour does not close the proof}

\begin{proposition}[Endpoint obstruction]\label{prop:contour}
At (13), the undeformed vertical contour in (15) is not absolutely
convergent for any real Fourier frequency $\alpha$.
\end{proposition}

\begin{proof}
For $y=i\lambda$, the two asymptotic formulas
\cite[equations (40)--(41)]{SS} give
\[
 \Gamma_M(y,m)\sim
 Z(m)e^{-\pi iB_{2,2}(y)/48},
\quad
 \Gamma_M(Q-y,-m)\sim
 Z(-m)^{-1}e^{+\pi iB_{2,2}(Q-y)/48}.
\tag{18}
\]
Since $B_{2,2}(Q-y)=B_{2,2}(y)$, the quadratic decay cancels and the
kernel tends to the nonzero constant $Z(m)/Z(-m)$.  The remaining
phase has modulus
\[
 \exp\left(-\frac{2\pi\alpha\lambda}{24\omega_1}\right).
\tag{19}
\]
It grows at one end when $\alpha\ne0$ and fails to decay at either end
when $\alpha=0$.
\end{proof}

There is nevertheless no finite pole pinch at $g=Q$.  The poles of the
first gamma factor lie in the cone at and to the left of $0$, while
those of the reflected factor lie in the cone at and to the right of
$g$.  At $g=Q>0$ the cones remain separated by $Q$.  Hence no
finite-pinch residue correction occurs; the obstruction is entirely at
imaginary infinity.

Likewise, the fused bilateral series has
$cd/(ab)=q^2$ and $z=-q$.  Slater's strict convergence annulus
$|q|^2<|q|<1$ collapses at the boundary.  Neither the original contour
nor the boundary bilateral sum proves Conjecture~\ref{conj:fusion};
equation (15) survives only as a meromorphic continuation.

\subsection*{Component type and the exact small divisor}

For finite frequency $(a,b)\bmod6$, choose the centered lift
\[
 s_{a,b}\equiv4b-5a\pmod6,\qquad
 s_{a,b}\in\{-3,-2,-1,0,1,2\}.
\]
Exactly six components have $s_{a,b}=0$; they are purely oscillatory
and admit the Fresnel/Abel prescription.  The remaining thirty grow at
one end of the vertical contour and require the interior strip and its
tilted finite part.  Reversing the convention for residue three changes
six growth orientations but not this $6+30$ split.

Put $\Delta(\tau)=\tau+\tau^{-1}-5$.  Direct calculation gives
\[
 A\tau-\tau
 =-\frac{24\tau}{24\tau-5}\Delta(\tau).
\tag{19a}
\]
Equivalently,
\[
 \gamma(s)=
 \frac{\beta+\beta^{-1}s^2+i\sqrt{21}s}{1+s^2},
\]
and therefore
\[
 \Delta(\gamma(s))
 =i\frac{21}{\beta}s+O(s^2),
\qquad
 2\pi|\Im(A\gamma(s)-\gamma(s))|
 =2\pi\sqrt{21}(1-\beta^{-6})s+O(s^2).
\tag{19b}
\]
This is the degenerating two-base decay coefficient.

The arithmetic at the same point is explicit:
\[
 \beta=[4;\overline{1,3}],\qquad
 (n\beta-m)(n\beta^{-1}-m)=m^2-5mn+n^2.
\]
Taking $m$ nearest to $n\beta$ and using the nonzero integral norm gives
\[
 \|n\beta\|\geq\frac1{\sqrt{21}\,n+\tfrac12},
\qquad
 |1-e^{2\pi in\beta}|
 \geq\frac4{\sqrt{21}\,n+\tfrac12}.
\tag{19c}
\]
Thus exponential damping of rate $\kappa(s)\asymp s$ competes with
quadratic-irrational divisors of order $1/n$ in the transition range
$n\asymp1/s$.  This is the sharp small-divisor form of MFC$_6$.

\begin{proposition}[Analytic stratum audit]\label{prop:stratum}
For both formal shifts, the six Fresnel frequency modes coincide
exactly with the six q-gamma singular-cancellation modes.  They do not
coincide with the conductor-lowered arithmetic stratum.
\end{proposition}

\begin{proof}
The Fresnel condition is $4b-5a\equiv0\pmod6$.  Inverting the two exact
TCC frequency maps sends its six solutions respectively to
\[
 \{(0,0),(4,1),(2,2),(0,3),(4,4),(2,5)\}
\]
and
\[
 \{(0,0),(5,5),(4,4),(3,3),(2,2),(1,1)\}.
\]
These are exactly the respective q-gamma cancellation patterns.  Each
set has denominator counts
\[
 (1,1,2,2)\quad\text{at denominators }(1,2,3,6).
\]
The thirty growing modes have counts $(0,2,6,22)$.  Only one point of
the proved three-point modulus-three orbit is Fresnel for either
shift.  The executable exact comparison is
\path{scripts/dimension_six_fresnel_stratum_audit.py}.
\end{proof}

\subsection*{A standalone sufficient estimate}

For $\tau\in\mathbb H$, put
\[
 \omega=24\tau-5,\quad q=e^{2\pi i\tau},\quad
 \widetilde q=e^{2\pi iA\tau},
\]
and define
\[
 G_\tau(\mu,m)=
 \frac{(\widetilde X(\mu,m);\widetilde q)_\infty}
 {(X(\mu,m);q)_\infty},
\]
where
\[
 X(\mu,m)=e^{2\pi i(\mu+m)/24},\quad
 \widetilde X(\mu,m)=\widetilde q\,
 e^{2\pi i(\mu+115m\omega)/(24\omega)}.
\]
With $D=4\tau-1$, let
\[
 \alpha_z=D(4b-5a)/3+2Dz,\qquad N_z=a+2-6z,
\]
\[
 K_{a,b}(z;\tau)=
 G_\tau(\alpha_z,N_z)G_\tau(-\alpha_z,4-N_z),
\]
\[
 \mathcal S_{a,b,r}(s)=
 \sum_{k\in\Z}K_{a,b}(r+3k;\gamma(s)),
 \qquad r=0,1,2.
\]
One alias period has multiplier $-q$; at the boundary the oscillations
are therefore $(-1)^ke^{2\pi ik\beta}$ and concentrate at
$|k|\asymp s^{-1}$.

\begin{conjecture}[Quantitative boundary fusion, BF$_6(\eta)$]
\label{conj:estimate}
For some $\eta,C,s_0>0$, uniformly over the thirty pairs
$4b-5a\not\equiv0\pmod6$ and $r=0,1,2$,
\[
 \left|\mathcal S_{a,b,r}(s)
 -\mathcal S^{\rm fus}_{a,b,r}\right|
 \le C\,|\Delta(\gamma(s))|^\eta,
 \qquad 0<s<s_0,
\tag{19d}
\]
where the fused value is the inherited $24$-factor Abel--Fresnel limit.
\end{conjecture}

A Dini modulus tending to zero would suffice instead of a power.
Conjecture~\ref{conj:estimate} implies MFC$_6$ and hence
Theorems~\ref{thm:stark} and~\ref{thm:conditional}.  The completely
standalone formulation is
\path{docs/dimension-six-standalone-estimate.md}.

\subsection*{Conditioning slopes}

After dividing the relative Arb radius by the requested
$q$-Pochhammer tolerance, the pinned asymptotic windows give
\[
\begin{array}{c|c|c|c|c}
d&\operatorname{Tr}(\beta)&\sqrt{\operatorname{disc}K}
&\dfrac{d\log_{10}C_d}{d(1/s)}
&\dfrac{d\log_{10}C_d}{d(1/|\tau-\beta|)}\\ \hline
4&3&\sqrt5&0.643602&1.43948\\
6&5&\sqrt{21}&2.803972&12.8721
\end{array}
\tag{19e}
\]
with $R^2=0.999933$ and $0.999996$, respectively, and replication at a
second precision.  This is essential-exponential conditioning in
$1/s$ for the factorized-continuation implementation.  It is not an
intrinsic exponent of BF$_6$.  The dimension-four value is proved
despite the same pathology; its slope is about $4.36$ times smaller.
See \path{scripts/dimension_six_conditioning_comparison.py}.

\section{Stabilizer ledger and conditional closure}
\label{sec:conditional}

\begin{lemma}[All multipliers]\label{lem:multipliers}
The matrix $A$ fixes every characteristic modulo $6$, satisfies
$\psi^2(A)=-1$, and for every
$\boldsymbol r=(a/6,b/6)$,
\[
 (\psi^{-2}\chi_{\boldsymbol r}^{-1})(A)
 =\Phi_{a,b}^{\,2}.
\tag{20}
\]
\end{lemma}

\begin{proof}
The first assertion is $A\equiv I\pmod6$.  The Rademacher invariant is
$6$, so
\[
 \psi^2(A)=e^{\pi i6/6}=-1.
\]
Substitution of $\boldsymbol r$ in Kopp's theta-character formula and
reduction modulo one agrees with the square of
\[
 \Phi_{a,b}
 =(-1)^{6+7(1+a)(1+b)}
 e^{-\pi i6/12}\tau_6^{-(a^2-5ab+b^2)}
\tag{21}
\]
in all $36$ cases.  This is exact rational arithmetic; see
\path{scripts/dimension_six_stabilizer_ledger.py}.
\end{proof}

\begin{proof}[Proof of Theorem~\ref{thm:conditional}]
Conjecture~\ref{conj:fusion} identifies the two-base boundary packet
with the AFK/Kopp packet characteristic by characteristic.
Lemma~\ref{lem:multipliers} shows that this identification has the
required stabilizer phase.  The exact finite frequency maps then turn
the two formal TCC sums into the two reconstructed matrices.  Their
trace is one and all $225$ rank-two minors vanish, so each has rank one
and is idempotent.  The matrix-to-TCC Fourier bridge gives both shifts
for the canonical tuple.  Integral covariance under
$\mathrm{GL}_2(\Z)$, applied also to the inverse transformation, gives
equality of shift sets and transports the result to every admissible
dimension-six form.
\end{proof}

\begin{proof}[Proof of Theorem~\ref{thm:stark}]
The exact characteristic-to-ray bridge identifies the fused boundary
packet with the differenced derivatives
$D_j=Z'_{(6)\infty_2}(0,g^j)$ in the certified Frobenius order.
Projecting to the primitive order-six Fourier line gives
\[
 L'_S(0,\chi_1)=r_0+\zeta_6r_1+\zeta_6^2r_2.
\tag{22}
\]
No TCC identity or minor vanishing is used in this projection.
\end{proof}

\subsection*{Grade-2 equivalence and Grade-3 attack surface}

The rigid endpoint content of Conjecture~\ref{conj:fusion} is
reduction-equivalent to (22).  Write
\[
 \Lambda_1=A_0+\zeta_6B_0,\qquad Q_3=D_0-D_1+D_2.
\]
The reverse reduction is
\[
 D_0=\frac{Q_3+2A_0+B_0}{3},\qquad
 D_1=\frac{-Q_3+A_0+2B_0}{3},\qquad
 D_2=\frac{Q_3-A_0+B_0}{3}.
\tag{23}
\]
Indeed, (22) fixes $\Lambda_1$, conjugation fixes $\Lambda_5$, and the
proved conductor-three identity fixes $Q_3$.  Equation (23), sign-class
reciprocity, shift/reflection/duplication, conductor lowering, and
Lemma~\ref{lem:multipliers} recover all thirty-six endpoint values.
Conversely, primitive Fourier projection gives (22).  Neither direction
uses TCC or a minor certificate.

The literal MFC$_6$ assertion is not derived from (22) by this standard
basis: equation (22) does not establish existence of the geodesic limit
or commutation of periodization with fusion.  Thus the endpoint value is
Grade-2 equivalent, but the converse for the full regularity hypothesis
is open.  Its family formulation has a genuinely
richer attack surface: one may differentiate in $\tau$, deform the
contour, track pinches and residue jumps, vary the lens label, iterate
the $A_6$ return map, and use badly-approximable-point estimates.  This
analytic structure is absent from the rigid equality (22).

\section{Scope and reproducibility}

The verified content of this note is the two-base reformulation,
the dimension-five/dimension-six dichotomy, the meromorphic
Sarkissian--Spiridonov specialization, the endpoint exclusion, and the
finite multiplier ledger.  The unconditional TCC is already proved in
dimensions four, five, seven, and eight in the companion
papers~\cite{PaperI,PaperII}.  Dimension six remains conditional on
Conjecture~\ref{conj:fusion}.

The cycle-by-cycle state and all negative-result gates are recorded in
\path{docs/dimension-six-state-notes-v3.md}.  The executable certificate
packet is regenerated by
\path{scripts/generate_dimension_six_amendment_certificates.sh}.
The zero-frequency tilted/Fresnel normalization gives the reciprocal
pair $-2\sqrt7\pm3\sqrt3$ and hence the independently enclosed trace
$-4\sqrt7$.  This is a normalization calibration only; it does not
imply the nonzero oriented fusion-continuity assertion.

\section*{Acknowledgments}

The author thanks the developers of PARI/GP, FLINT, Arb, and
python-flint.  The formulation in this note benefited from repeated
adversarial checks of normalization, place labeling, and boundary
convergence.

\begin{thebibliography}{9}

\bibitem{AFK}
D.~M. Appleby, S.~T. Flammia, and G.~S. Kopp,
\emph{A constructive approach to Zauner's conjecture via the
Shintani--Faddeev modular cocycle},
arXiv:2501.03970.

\bibitem{Kopp}
G.~S. Kopp,
\emph{The Shintani--Faddeev modular cocycle: Stark units from
$q$-Pochhammer ratios},
arXiv:2411.06763.

\bibitem{SS}
G.~A. Sarkissian and V.~P. Spiridonov,
\emph{General modular quantum dilogarithm and beta integrals},
arXiv:1910.11747, version 4.

\bibitem{PaperI}
H.~Zhao,
\emph{Twisted-convolution identities in dimensions four and five
from Shintani ray units}, preprint, 2026.

\bibitem{PaperII}
H.~Zhao,
\emph{Twisted-convolution identities in dimensions seven and eight:
Shintani height rigidity and CM descent}, preprint, 2026.

\end{thebibliography}

\end{document}
